\chapter{Difference D'Artagnan Finds Between the Intendant and the Superintendent}
\chaptermark{Difference Between the Intendant and the Superintendent}


\lettrine{M}{.} Colbert resided in the Rue Neuve des Petits-Champs, in  a house which had belonged to Beautru. D'Artagnan's legs cleared  the distance in a short quarter of an hour. When he arrived at the residence of  the new favourite, the court was full of archers and police, who came to  congratulate him, or to excuse themselves, according to whether he should  choose to praise or blame. The sentiment of flattery is instinctive with people  of abject condition; they have the sense of it, as the wild animal has that of  hearing and smell. These people, or their leader, understood that there was a  pleasure to offer to M. Colbert, in rendering him an account of the fashion in  which his name had been pronounced during the rash enterprise of the morning.  D'Artagnan made his appearance just as the chief of the watch was giving  his report. He stood close to the door, behind the archers. That officer took  Colbert on one side, in spite of his resistance and the contradiction of his  bushy eyebrows. <In case,> said he, <you really desired,  monsieur, that the people should do justice on the two traitors, it would have  been wise to warn us of it; for, indeed, monsieur, in spite of our regret at  displeasing you, or thwarting your views, we had our orders to execute.>

<Triple fool!> replied Colbert, furiously shaking his hair, thick  and black as a mane; <what are you telling me? What! that I could have  had an idea of a riot! Are you mad or drunk?>

<But, monsieur, they cried <Vive Colbert!>> replied the  trembling watch.

<A handful of conspirators\longdash>

<No, no; a mass of people.>

<Ah! indeed,> said Colbert, expanding. <A mass of people  cried <Vive Colbert!> Are you certain of what you say,  monsieur?>

<We had nothing to do but open our ears, or rather to close them, so  terrible were the cries.>

<And this was from the people, the real people?>

<Certainly, monsieur; only these real people beat us.>

<Oh! very well,> continued Colbert, thoughtfully. <Then you  suppose it was the people alone who wished to burn the condemned?>

<Oh! yes, monsieur.>

<That is quite another thing. You strongly resisted, then?>

<We had three of our men crushed to death, monsieur!>

<But you killed nobody yourselves?>

<Monsieur, a few of the rioters were left upon the square, and one among  them who was not a common man.>

<Who was he?>

<A certain Menneville, upon whom the police have a long time had an  eye.>

<Menneville!> cried Colbert, <what, he who killed Rue de la  Huchette, a worthy man who wanted a fat fowl?>

<Yes, monsieur; the same.>

<And did this Menneville also cry, <Vive Colbert>?>

<Louder than all the rest; like a madman.>

Colbert's brow grew dark and wrinkled. A kind of ambitious glory which  had lighted his face was extinguished, like the light of glow-worms we crush  beneath the grass. <Then you say,> resumed the deceived intendant,  <that the initiative came from the people? Menneville was my enemy; I  would have had him hung, and he knew it well. Menneville belonged to the Abbé  Fouquet—the affair originated with Fouquet; does not everybody know that  the condemned were his friends from childhood?>

<That is true,> thought D'Artagnan, <and thus are all  my doubts cleared up. I repeat it, Monsieur Fouquet may be called what they  please, but he is a very gentlemanly man.>

<And,> continued Colbert, <are you quite sure Menneville is  dead?>

D'Artagnan thought the time was come for him to make his appearance.  <Perfectly, monsieur;> replied he, advancing suddenly.

<Oh! is that you, monsieur?> said Colbert.

<In person,> replied the musketeer with his deliberate tone;  <it appears that you had in Menneville a pretty enemy.>

<It was not I, monsieur, who had an enemy,> replied Colbert;  <it was the king.>

<Double brute!> thought D'Artagnan, <to think to play  the great man and the hypocrite with me. Well,> continued he to Colbert,  <I am very happy to have rendered so good a service to the king; will you  take upon you to tell his majesty, monsieur l'intendant?>

<What commission is this you give me, and what do you charge me to tell  his majesty, monsieur? Be precise, if you please,> said Colbert, in a  sharp voice, tuned beforehand to hostility.

<I give you no commission,> replied D'Artagnan, with that  calmness which never abandons the banterer; <I thought it would be easy  for you to announce to his majesty that it was I who, being there by chance,  did justice upon Menneville and restored order to things.>

Colbert opened his eyes and interrogated the chief of the watch with a  look—<Ah! it is very true,> said the latter, <that this  gentleman saved us.>

<Why did you not tell me, monsieur, that you came to relate me  this?> said Colbert with envy; <everything is explained, and more  favourably for you than for anybody else.>

<You are in error, monsieur l'intendant, I did not at all come for  the purpose of relating that to you.>

<It is an exploit, nevertheless.>

<Oh!> said the musketeer carelessly, <constant habit blunts  the mind.>

<To what do I owe the honour of your visit, then?>

<Simply to this: the king ordered me to come to you.>

<Ah!> said Colbert, recovering himself when he saw D'Artagnan  draw a paper from his pocket; <it is to demand some money of me?>

<Precisely, monsieur.>

<Have the goodness to wait, if you please, monsieur, till I have  dispatched the report of the watch.>

D'Artagnan turned upon his heel, insolently enough, and finding himself  face to face with Colbert, after his first turn, he bowed to him as a harlequin  would have done; then, after a second evolution, he directed his steps towards  the door in quick time. Colbert was struck with this pointed rudeness, to which  he was not accustomed. In general, men of the sword, when they came to his  office, had such a want of money, that though their feet seemed to take root in  the marble, they hardly lost their patience. Was D'Artagnan going  straight to the king? Would he go and describe his rough reception, or recount  his exploit? This was a matter for grave consideration. At all events, the  moment was badly chosen to send D'Artagnan away, whether he came from the  king, or on his own account. The musketeer had rendered too great a service,  and that too recently, for it to be already forgotten. Therefore Colbert  thought it would be better to shake off his arrogance and call D'Artagnan  back. <Ho! Monsieur d'Artagnan,> cried Colbert, <what!  are you leaving me thus?>

D'Artagnan turned round: <Why not?> said he, quietly,  <we have no more to say to each other, have we?>

<You have, at least, money to receive, as you have an order?>

<Who, I? Oh! not at all, my dear Monsieur Colbert.>

<But, monsieur, you have an order. And, in the same manner as you give a  sword-thrust, when you are required, I, on my part, pay when an order is  presented to me. Present yours.>

<It is useless, my dear Monsieur Colbert,> said D'Artagnan,  who inwardly enjoyed this confusion in the ideas of Colbert; <my order is  paid.>

<Paid, by whom?>

<By monsieur le surintendant.>

Colbert grew pale.

<Explain yourself,> said he, in a stifled voice—<if you  are paid why do you show me that paper?>

<In consequence of the word of order of which you spoke to me so  ingeniously just now, dear M. Colbert; the king told me to take a quarter of  the pension he is pleased to make me.>

<Of me?> said Colbert.

<Not exactly. The king said to me: <Go to M. Fouquet; the  superintendent will, perhaps, have no money, then you will go and draw it of M.  Colbert.>>

The countenance of M. Colbert brightened for a moment; but it was with his  unfortunate physiognomy as with a stormy sky, sometimes radiant, sometimes dark  as night, according as the lightening gleams or the cloud passes. <Eh!  and was there any money in the superintendent's coffers?> asked he.

<Why, yes, he could not be badly off for money,> replied  D'Artagnan—<it may be believed, since M. Fouquet, instead of  paying me a quarter or five thousand livres\longdash>

<A quarter or five thousand livres!> cried Colbert, struck, as  Fouquet had been, with the generosity of the sum for a soldier's pension,  <why, that would be a pension of twenty thousand livres?>

<Exactly, M. Colbert. Peste! you reckon like old Pythagoras; yes, twenty  thousand livres.>

<Ten times the appointment of an intendant of the finances. I beg to  offer you my compliments,> said Colbert, with a vicious smile.

<Oh!> said D'Artagnan, <the king apologized for giving  me so little; but he promised to make it more hereafter, when he should be  rich; but I must be gone, having much to do\longdash>

<So, then, notwithstanding the expectation of the king, the  superintendent paid you, did he?>

<In the same manner, as, in opposition to the king's expectation,  you refused to pay me.>

<I did not refuse, monsieur, I only begged you to wait. And you say that  M. Fouquet paid you your five thousand livres?>

<Yes, as you might have done; but he did even better than that, M.  Colbert.>

<And what did he do?>

<He politely counted me down the sum-total, saying, that for the king,  his coffers were always full.>

<The sum-total! M. Fouquet has given you twenty thousand livres instead  of five thousand?>

<Yes, monsieur.>

<And what for?>

<In order to spare me three visits to the money-chest of the  superintendent, so that I have the twenty thousand livres in my pocket in good  new coin. You see, then, that I am able to go away without standing in need of  you, having come here only for form's sake.> And D'Artagnan  slapped his hand upon his pocket, with a laugh which disclosed to Colbert  thirty-two magnificent teeth, as white as teeth of twenty-five years old, and  which seemed to say in their language: <Serve up to us thirty-two little  Colberts, and we will chew them willingly.> The serpent is as brave as  the lion, the hawk as courageous as the eagle, that cannot be contested. It can  only be said of animals that are decidedly cowardly, and are so called, that  they will be brave only when they have to defend themselves. Colbert was not  frightened at the thirty-two teeth of D'Artagnan. He recovered, and  suddenly,—<Monsieur,> said he, <monsieur le  surintendant has done what he had no right to do.>

<What do you mean by that?> replied D'Artagnan.

<I mean that your note—will you let me see your note, if you  please?>

<Very willingly; here it is.>

Colbert seized the paper with an eagerness which the musketeer did not remark  without uneasiness, and particularly without a certain degree of regret at  having trusted him with it. <Well, monsieur, the royal order says  thus:—<At sight, I command that there be paid to M.  d'Artagnan the sum of five thousand livres, forming a quarter of the  pension I have made him.>>

<So, in fact, it is written,> said D'Artagnan, affecting  calmness.

<Very well; the king only owed you five thousand livres; why has more  been given to you?>

<Because there was more; and M. Fouquet was willing to give me more; that  does not concern anybody.>

<It is natural,> said Colbert with a proud ease, <that you  should be ignorant of the usages of state-finance; but, monsieur, when you have  a thousand livres to pay, what do you do?>

<I never have a thousand livres to pay,> replied D'Artagnan.

<Once more,> said Colbert, irritated—<once more, if you  had any sum to pay, would you not pay what you ought?>

<That only proves one thing,> said D'Artagnan; <and  that is, that you have your own particular customs in finance, and M. Fouquet  has his own.>

<Mine, monsieur, are the correct ones.>

<I do not say that they are not.>

<And you have accepted what was not due to you.>

D'Artagnan's eyes flashed. <What is not due to me yet, you  meant to say, M. Colbert; for if I have received what was not due to me at all,  I should have committed a theft.>

Colbert made no reply to this subtlety. <You then owe fifteen thousand  livres to the public chest,> said he, carried away by his jealous ardour.

<Then you must give me credit for them,> replied D'Artagnan,  with his imperceptible irony.

<Not at all, monsieur.>

<Well! what will you do, then? You will not take my rouleaux from me,  will you?>

<You must return them to my chest.>

<I! Oh! Monsieur Colbert, don't reckon upon that.>

<The king wants his money, monsieur.>

<And I, monsieur, I want the king's money.>

<That may be so; but you must return this.>

<Not a sou. I have always understood that in matters of comptabilite, as  you call it, a good cashier never gives back or takes back.>

<Then, monsieur, we shall see what the king will say about it. I will  show him this note, which proves that M. Fouquet not only pays what he does not  owe, but that he does not even take care of vouchers for the sums that he has  paid.>

<Ah! now I understand why you have taken that paper, M. Colbert!>

Colbert did not perceive all that there was of a threatening character in his  name pronounced in a certain manner. <You shall see hereafter what use I  will make of it,> said he, holding up the paper in his fingers.

<Oh!> said D'Artagnan, snatching the paper from him with a  rapid movement; <I understand perfectly well, M. Colbert; I have no  occasion to wait for that.> And he crumpled up the paper he had so  cleverly seized.

<Monsieur, monsieur!> cried Colbert, <this is  violence!>

<Nonsense! You must not be particular about a soldier's  manners!> replied D'Artagnan. <I kiss your hands, my dear M.  Colbert.> And he went out, laughing in the face of the future minister.

<That man, now,> muttered he, <was about to grow quite  friendly; it is a great pity I was obliged to cut his company so soon.>  